\section{Discussion}
\label{sec:discussion}

\paragraph{What the solver loop buys.} The loop turns a model's claim (``this puzzle has one answer'') into a certificate, but the certificate covers the encoding, not the text. Our results suggest three practical rules. First, the translation from text to constraints is reliable when the text is regular: on template text no encoder made any error. Second, dual encoding catches independent misreadings, which are what a second model run is good at, but not errors that two runs share, and two runs of the same model share a great deal: the two Sonnet encoders agreed on every paraphrased puzzle, including the seven they both got wrong. Third, uniqueness is a strong and cheap second check for minimal puzzles, because a minimal puzzle has no slack: a misread clue almost always removes the intended solution or lets in another. The combination caught every wrong encoding we observed. What it cannot catch is an error that leaves the puzzle unique with the same answer, such as our own ambiguous template; such errors do not change the answer but do mean the text says something other than the certificate, which matters if the puzzle is used to teach or to explain.

\paragraph{LLM-written puzzles.} Sonnet's puzzles were valid far more often than we expected, and one round of solver feedback fixed the rest. The problem is not validity but quality: they are easy, padded with redundant clues, and in Sudoku they reuse a famous layout. A solver that minimises and grades the puzzle fixes the first two; the third is a warning for anyone who builds evaluation sets by asking a model for puzzles.

\paragraph{The ladder.} Five nested propagation engines give a level that is cheap to compute, interpretable in solving terms, and closely related to solver effort and to the size of minimal explanations in every family. It does not, at these sizes, predict whether current LLMs succeed, because they almost always do. The exception is the smallest model on $9{\times}9$ Sudoku, where failures, and especially give-ups, concentrate in the higher levels, though our data cannot separate level from batch effects there. A ladder whose top levels are rare in random generation (no Sudoku above L3 in 9{,}500) is also a ladder whose top levels must be searched for, which is what E4 attempts.

\paragraph{Fair but hard.} We did not find fair zebra puzzles that Haiku fails. This is a null result at our scale (two runs, at most three generations, 8 individuals) and it does not show that such puzzles do not exist; but the space we searched, all unique L1--L3 puzzles of size up to $6{\times}6$, was solved without error in \EFourZebraEvals of \EFourZebraEvals attempts. Larger instances, other genres, or fitness based on solving cost rather than errors may behave differently.

\section{Limitations}
\label{sec:limitations}

\begin{itemize}[leftmargin=1.2em,itemsep=1pt,topsep=2pt]
\item \emph{How the models were called.} No API key was available, so every call was a fresh Claude Code subagent with no shared context (Section~\ref{sec:llmcalls}). Temperature and reasoning settings were not under our control, model versions are identified by alias, and the workers had a general system prompt we did not write. The audit of tool calls (\NAudited of \NCalls calls matched to transcripts) found no code execution; \NForbidden calls made one extra status or messaging call with no puzzle content (one edited its own output file), and \NExtraWrite wrote their output in several steps. Calls with extra tool use were kept.
\item \emph{Scale.} The study is small: 299 puzzles in E1, 100 in E2, 116 in E3, three evolution runs. Rates near zero have wide intervals, and E3 has 8 puzzles per cell (4 for Killer L5).
\item \emph{Batching.} Several items share a call. Items in a call are not independent, and a call that gives up affects all its items (one Haiku Sudoku call left five items unanswered, which we score as wrong).
\item \emph{Augmented L1 puzzles.} L1 cells of three families were filled by adding redundant clues to minimal puzzles, so they differ from the other cells in clue count and redundancy.
\item \emph{Ground truth in E1} comes from our generator; one of its templates was ambiguous, which we report as a finding rather than an encoder error.
\item \emph{Cost} is estimated from total token counts with a conservative split between input and output and list prices without caching; real costs would be lower.
\item \emph{Ladder.} The ladder is defined on one CNF encoding; a different encoding (for example with auxiliary variables) would change what unit propagation can do, and humans' notion of difficulty is not modelled.
\end{itemize}
